% Source: https://artofproblemsolving.com/community/c3h1741490_ratio_of_height_and_width
% Stable identity: aops-topic-1741490-q01
Three circles are inscribed in a rectangle of width $w$ and height $h$, as shown. Two of the circles are congruent; one is tangent to the left and top sides, the other to the left and bottom sides, and they are tangent to each other. The third circle is larger and is tangent to the top, bottom and right sides of the rectangle and to both smaller circles. What is the ratio of $h$ to $w$? Express your answer as a decimal to the nearest hundredth.

\begin{asy}
usepackage("amsmath");
size(5cm);
defaultpen(fontsize(10pt));
real r=1,w=3+2*sqrt(2),h=4;
draw((0,0)--(w,0)--(w,h)--(0,h)--cycle,linewidth(0.9));
draw(circle((1,1),1));draw(circle((1,3),1));draw(circle((w-2,2),2));
draw((0,-0.45)--(w,-0.45),Arrows(4));label("$w$",(w/2,-0.45),S);
draw((w+0.45,0)--(w+0.45,h),Arrows(4));label("$h$",(w+0.45,h/2),E);
label("\mbox{}",(min(currentpicture,user=true)+max(currentpicture,user=true))/2,black); // keeps the text after the figure black when LaTeX typesets the labels inline
\end{asy}
